\section{Introduction}\label{sec:intro}
Locomotion policies for legged robots are commonly learned in simulation and transferred to hardware \cite{hwangbo2019agile}, with increasingly detailed actuator models, yet these models treat the electrical supply of each joint as ideal or fixed. The energy model of Bjelonic et al.\ does not include battery dynamics \cite{bjelonic2026pace}, and the actuator models of Tao et al.\ \cite{tao2026actuator} and Zhang et al.\ \cite{zhang2026vra} bound each motor's voltage by a given direct-current (DC) bus voltage. On a battery-powered robot, however, all drives draw from one battery through one DC bus. The bus voltage falls with the total load, and the battery current has a limit. A lower bus voltage in turn reduces the torque that every joint can produce at speed.

Hardware work points to this gap; energy-aware work targets a different quantity. Miller et al.\ attribute a difference between simulated and real swing-leg motion to power distribution limits that their simulation does not include \cite{miller2025spot}. Nai et al.\ note that common simulators model battery consumption inaccurately or not at all, and learn it from real data \cite{nai2025power}. A battery model with transient voltage sag has been built for a humanoid, with its coupling to robot control left to future work \cite{evagora2025battery}. Energy-aware methods reduce the energy consumed over a motion or a run \cite{valsecchi2024power,huang2026eco,lee2026raibo2}. This paper concerns the instantaneous limits of the supply instead: a lowest admissible bus voltage and a largest battery current.

We model one leg of a quadruped the size of the Unitree Go1, with two drives on a shared battery and DC bus, and enforce these limits with a filter between a learned policy and the drives. At every simulation substep the filter turns two discrete-time barrier conditions, on the bus voltage and on the battery current, into one power budget for the next substep, and projects the drive currents onto it. We train policies with soft actor-critic (SAC) with the filter in the loop (filter-trained policies). We compare them across states of charge, parameter variants and errors in the battery model with policies trained without it (baseline policies), with the same filter added only at test time (test-time filter), with reward penalties and with a Lagrangian constrained method. The artificial-intelligence contribution is thus a physics-derived safety filter inside the reinforcement-learning training loop, with a one-step guarantee under stated assumptions whenever currents meeting the budget exist. The contributions are:
\begin{enumerate}
\item an electrically coupled simulation layer: a battery fitted to public cell data, a DC link with brake chopper and averaged drives with current and voltage limits, checked against a switching-level drive model (Section~\ref{sec:model});
\item a barrier filter for bus undervoltage and battery overcurrent. A load meets both barrier conditions in the next substep if and only if it does not exceed a power budget computed in closed form (Proposition~\ref{prop:budget}); the filter projects the drive currents onto this budget by bisection on one multiplier, with a one-step guarantee under stated assumptions in every substep in which currents meeting the budget exist (Corollary~\ref{cor:closed});
\item a comparison of the filter-trained policies with the test-time filter, reward penalties and a Lagrangian method (Section~\ref{sec:results}), and an ablation study over initialisation, filter placement, barrier rate, margins and a joint-level limit (Section~\ref{sec:ablation}). Penalty policies with the test-time filter met the limits and hopped higher than the filter-trained policies in most settings, but one of their seeds fell in stress tests in which the filter-trained policies did not (Section~\ref{sec:testtime});
\item a robustness study over parameter variants, barrier rates and errors in the battery's electromotive force and resistance, with a first-order sensitivity rule for the current margin (Remark~\ref{rem:sens}) and an extended Kalman filter (EKF) that supplies the battery parameters, including the conditions under which it fails; margins scaled by the EKF's confidence were tested as an option and acted as a switch (Sections~\ref{sec:estimator} and~\ref{sec:estresults}).
\end{enumerate}
The budget differs from a fixed or state-dependent power cap set by hand in two ways. First, it is the exact one-step admissible power of the bus model: a load meets both barrier conditions in the next substep if and only if it does not exceed the budget (Proposition~\ref{prop:budget}), which is recomputed at every substep from the bus voltage, the battery current and the battery parameters. Second, both drives share one budget, so one joint can use power that the other does not need. Splitting the same budget equally between the two drives, a joint-level limit, acts more often and gives lower hops (Section~\ref{sec:pjm}). For the filter-trained policies it is active in \res{pm_perjoint_filtws_pjtest_filt}\% of the substeps against \res{pm_base_filtws_filt}\% for the shared budget, and their max height at state of charge 0.2 is \res{abl_perjoint_filtws_pjtest_soc0.2_hmax} against \res{abl_base_filtws_soc0.2_hmax}~m.

Figure~\ref{fig:framework} gives an overview of the framework and maps its parts to these contributions.
\begin{figure}[tbp]\centering\genfig{F0_framework}{\linewidth}\caption{Overall framework: (a) public sources and assumed values behind the model parameters and (b) the electrically coupled plant (Section~\ref{sec:model}), (c) learning with the barrier power filter (Sections~\ref{sec:filter} and~\ref{sec:setup}) and (d) evaluation (Sections~\ref{sec:results} and~\ref{sec:ablation}). The strip maps the panels to the contributions.}\label{fig:framework}\end{figure}
